\section{결과}\label{sec:results-ko}

\subsection{평지 절제 실험}

\begin{figure*}[t]
\centering\includegraphics[width=0.91\textwidth]{flat_metrics.pdf}
\caption{동일 모델을 이용한 평지 시험 결과(제어기당 1회, 측정 시간 6초).
조건별 단일 시행의 기술 통계이므로 오차 막대는 표시하지 않았다.}
\label{fig:flat-ko}
\end{figure*}

표~\ref{tab:flat-ko}와 그림~\ref{fig:flat-ko}는 세 제어 방식의 서로 다른
장단점을 보여준다. 협응 구름 주행의 평균 속도는
$0.200\,\mathrm{m\,s^{-1}}$로, $0.18\,\mathrm{m\,s^{-1}}$ 명령보다 $11.2\%$
높고 관절 보행보다 $71\%$ 빨랐다. 이 결과는 기술적으로 H1과 일치하지만, 보행
파라미터가 서로 달라 인과 관계를 주장할 수는 없다. 더 강한 해석, 즉
협응 구름 주행이 명령 추종 성능까지 가장 좋다는 해석은 같은 결과로 반박된다.
일정한 명령에 대한 속도 추종 오차는 관절 보행에서 가장 작았다
($e_v=0.097$ 대 $0.253$, $0.268\,\mathrm{m\,s^{-1}}$). 구름 주행은 명령
속도를 정확히 추종한 것이 아니라 초과했기 때문이다. 또한 전체 구름 주행은
관절 보행보다 절대 기계일이 $3.18\times$, 접촉 미끄럼이 $6.23\times$ 컸다.
말단 단독 회전은 관절 보행보다 느리고 에너지 비용도 컸으므로, 호의 회전만으로는
속도 향상을 설명할 수 없다. 이 모델에서는 몸체 자세와 말단 회전을 함께 제어한
조건이 가장 빨랐으며, 그 대가로 기계일과 미끄럼이 증가했다.

\begin{table}[t]
\caption{0.18 m/s 전진 명령 평지 공칭 시행}
\label{tab:flat-ko}
\centering\small
\begin{tabular}{lrrrr}
\toprule
제어기 & $\bar v_x$ & $e_v$ & $C_{\rm mt}$ & 미끄럼\\
 & m/s & m/s & -- & m\\
\midrule
관절 보행 & 0.117 & \textbf{0.097} & \textbf{0.661} & \textbf{0.212}\\
말단 단독 & 0.090 & 0.268 & 2.576 & 1.603\\
협응 호 이동 & \textbf{0.200} & 0.253 & 1.227 & 1.323\\
\bottomrule
\end{tabular}
\end{table}

측정된 미끄럼은 원인을 알 수 없는 시뮬레이션 인공물이 아니다.
식~(\ref{eq:contactspeed-ko})에 따르면 $\alpha=1$일 때 지지기 접촉면의 명령
속도는 $v_c=0.411\,\mathrm{m\,s^{-1}}$로, 구름 조건에서 달성한 몸체 속도의 약
두 배다. 식~(\ref{eq:slippredict-ko})은 6초 동안 누적되는 미끄럼을
$|v_c-\bar v_x|T=1.27\,$m로 예측하며, 측정값 $1.323\,$m와의 오차는 $4.2\%$다.
말단 단독 회전 조건에서는 같은 식이 $1.93\,$m를 예측했지만 측정값은
$1.603\,$m로, 예측값이 $20\%$ 더 컸다. 몸체가 접촉 패치의 진행 방향을 충분히
따라가지 못하면 식~(\ref{eq:align-ko})의 정렬도 $A_i$가 1보다 작아지기 때문이다.
이 계산은 독립적인 예측이 아니라 일관성 검산이다. $\bar v_x$ 자체가 같은 시행에서
측정되었고 접촉 정렬도도 주기 중 변하기 때문이다. 그럼에도 고정된 활성도 기반
회전 명령이 접촉면 속도와 몸체 속도 사이에 큰 차이를 만들 수 있음을 보여준다.
따라서 보고한 미끄럼은 호 형상의 고유 성질이 아니라 현재 제어 설정의 특성이다.

\subsection{계단 등반 협응 절제 실험}

\begin{figure}[t]
\centering\includegraphics[width=\columnwidth]{stair_results.pdf}
\caption{세 가지 제어 제한 조건과 세 종류의 3단 계단에서 측정한 상단 도달 시간
(조건당 결정론적 1회). \emph{no top}은 상단에 도달하지 못한 경우를 뜻하며,
성공 여부는 표~\ref{tab:stairs-ko}에 별도로 정리했다.}
\label{fig:stairs-ko}
\end{figure}

표~\ref{tab:stairs-ko}는 동일한 개방형 호 모델에서 세 제어 제한 조건을 비교한다.
$100\,$mm 계단에서는 모든 전략이 성공했다. $150\,$mm$(1.22r_o)$ 계단에서는
관절 운동을 전혀 사용하지 않은 말단 단독 회전도 $14.960\,$s 만에 상단에
도달했다. 이 단높이는 개구 현보다 짧으므로($150<l_o=226\,$mm), 결과는 계단
모서리의 개구부 진입과 일치한다. 그러나 축토크를 가하는 닫힌 바퀴도 계단을 오를
수 있으므로 이 결과만으로 개구부 효과를 분리할 수 없다. 이를 위해서는
\ref{sec:method-ko}절의 동일 조건 닫힌 바퀴 대조군이 필요하다.
$200\,$mm$(1.63r_o)$ 계단에서는 두 제한 조건이 모두 실패했고, 협응 제어기만
$5.964\,$s, $88.18\,$J로 상단에 도달했다. 이는 시험한 개방형 호 제어 제한
내에서만 H2를 지지한다. 한편 등반 속도와 안정성은 반드시 일치하지
않았다. $150\,$mm 계단에서 개루프 제어기는 $5.546\,$s로 가장 빨랐지만, 직립도는
$u_z=0.459$까지 낮아져 $\cos60^\circ$의 종료 경계에 근접했다. 반면 협응
제어기는 $u_z=0.755$를 유지했다.

\begin{table}[t]
\caption{계단 상단 도달 결과(결정론적 1회)}
\label{tab:stairs-ko}
\centering\small
\resizebox{\columnwidth}{!}{%
\begin{tabular}{lrrrrrr}
\toprule
&\multicolumn{2}{c}{100 mm 단높이}&\multicolumn{2}{c}{150 mm 단높이}&\multicolumn{2}{c}{200 mm 단높이}\\
제어기 & 도달 & s & 도달 & s & 도달 & s\\
\midrule
말단 단독 & 예 & 4.886 & 예 & 14.960 & 아니오 & --\\
개루프 & 예 & 3.818 & 예 & 5.546 & 아니오 & --\\
전체 제어기 & 예 & 2.646 & 예 & 6.736 & \textbf{예} & 5.964\\
\bottomrule
\end{tabular}}
\end{table}

\subsection{모드 전환 및 불규칙 지형 기초 시험}

보행에서 구름 주행으로의 전환 조건은 $0.342\,$s, 구름 주행에서 보행으로의 전환
조건은 $0.660\,$s 만에 충족되었다. 전환 및 회복 구간 전체에서 측정한 최소
직립도는 두 방향 모두 $0.991$이었다. 금지된 충돌이나 역기구학 실패도 발생하지
않았으므로 어느 전환에서도 H3의 기울기 한계를 위반하지 않았다. 무작위 불규칙
지형에서는 총 9회의 짧은 시행이 모두 기울기 한계로 종료되지 않고 완료되었다.
3개 난수 시드에서 얻은 평균 $\bar v_x$는 관절 보행, 말단 단독 회전, 협응 호
주행 순으로 $0.0768\pm0.0043$, $0.1025\pm0.0356$,
$0.2009\pm0.0126\,\mathrm{m\,s^{-1}}$였다. 그러나 시드가 3개뿐이고 측정 시간도
$2\,$s이므로, 이 결과만으로 불규칙 지형에서 특정 제어 방식이 더 우수하다고
결론 내릴 수는 없다.
